\chapter{Расчёт}

\section{Задание}
Помещение для программистов имеет размеры: длину А = 12 м; ширину В = 8 м; высоту Н = 4 м. Высота рабочей поверхности $h_{\text{рп}}$ = 0,8 м. Требуется рассчитать общее равномерное освещение помещения при использовании светильников типа ЛОУ 1П с двумя люминесцентными лампами и составить схему размещения светильников в помещении.

\section{Решение}

Определяем для помещения норму освещённости согласно таблице 2.1 $E = 300$ лк.

Выбираем из табл. 2.2 с учетом состояния воздушной среды в помещении для программистов коэффициент запаса $K_{\text{з}} = 1,5$.

Принимаем расстояние от потолка до светильника $h_{\text{св}}$ = 0,5 м, исходя из рекомендуемых значений $h_{\text{св}} \le 2 $ м.

Определяем высоту установки светильников над рабочей поверхностью h: 

\[
h =  H - ( h_{\text{св}} + h_{\text{рп}} ) = 4 - (0,5 + 0,8) = 2,7 \text{м}.
\]

Из светотехнических данных светильника ЛСП 01 (см. табл. 2.3) выбираем наивыгоднейшее отношение $\lambda_\text{нв}=\frac{L}{h}=1,5$

Определяем расстояние между соседними рядами светильников L:

\[
L=\lambda_\text{нв} \cdot h=1,5 \cdot 2,7 = 4,05 \text{м}
\]

Тогда при ширине помещения В = 8 м число рядов светильников $n_{B}$ будет

\[
n_{B} = \frac{B}{L}=\frac{8}{4,05}=1,975
\]

Принимаем число рядов светильников в помещении $n_{B} = 2$.

Для определения величины коэффициента использования светового потока $\eta$ выбираем из табл. 2.4 коэффициенты отражения $\rho$ поверхностей помещения для программистов: потолка $\rho_\text{п} = 70\% $; стен $\rho_\text{ст} = 50\%$; рабочей поверхности $\rho_\text{рп} = 10 \%$ и рассчитываем индекс помещения $i$ по формуле:

\[
i=\frac{A \cdot B}{h \cdot (A+B)}=\frac{12*8}{2,7*(12+8)}=1,77
\]

Значение индекса помещения принимается равным ближайшему значению в табл. 2.5. В рассчитываемом примере принимаем $i = 1,75$.

Для найденных значений $i$, $\rho_\text{п}$, $\rho_\text{ст}$, $\rho_\text{рп}$ и светильника ЛОУ 1П находим из табл. 2.5 коэффициент использования светового потока $\eta = 0,53$.

Определяем световой поток светильников одного ряда $\text{Ф}_\text{р}$, подставив в формулу (2.1) вместо количества светильников $N$ число их рядов $n_\text{в}$:

\[
\text{Ф}_\text{р} = \frac{E \cdot K_\text{з} \cdot S \cdot z}{N \cdot \eta} = \frac{300 \cdot 1,5 \cdot 96 \cdot 1,1 }{2 \cdot 0,53} = \frac{47520}{1,06}=44830 \text{ лм}
\]

Рассчитаем необходимое количество в ряду двухламповых светильников $n_A$ для создания светового потока одного ряда $\text{Ф}_\text{р}$:

\[
n_A = \frac{\text{Ф}_\text{р}}{2 \cdot \text{Ф}_\text{л}}
\]

где $\text{Ф}_\text{л}$ – световой поток стандартной лампы в светильнике, лм.
Так как согласно данным табл. 2.3 в светильнике ЛОУ 1П применяются две люминесцентные лампы мощностью 40 Вт или две лампы мощностью 80 Вт, то для осветительной установки можно выбрать одну из следующих стандартных ламп с соответствующим световым потоком $\text{Ф}_\text{л}$ (см. табл. 2.6):

\begin{table}
	\caption{}
	\centering
	\begin{tabular}{|c|c|c|c|}
		\hline
		Лампа & $\text{Ф}_\text{л}$ & $n_A$ & $L_\text{р} = L_\text{св} \cdot n_A$ \\ \hline
		ЛДЦ 40 & 2200 & $\frac{44830}{2 \cdot 2200} \approx 10$ & $1,24 \cdot 10 = 12,4$ м \\ \hline
		ЛД 40 & 2500  & $\frac{44830}{2 \cdot 2500} \approx 9$ & $1,24 \cdot 9 = 11,16$ м \\ \hline
		ЛХБ 40 & 3000  & $\frac{44830}{2 \cdot 3000} \approx 8$ & $1,24 \cdot 8 = 9,92$ м \\ \hline
		ЛТБ 40 & 3100  & $\frac{44830}{2 \cdot 3100} \approx 8$ & $1,24 \cdot 8 = 9,92$ м \\ \hline
		ЛБ 40 & 3200  & $\frac{44830}{2 \cdot 3200} \approx 8$ & $1,24 \cdot 8 = 9,92$ м \\ \hline
		ЛДЦ 80 & 3800  & $\frac{44830}{2 \cdot 3800} \approx 6$ & $1,536 \cdot 6 = 9,216$ м \\ \hline
		ЛД 80 & 4300  & $\frac{44830}{2 \cdot 4300} \approx 6$ & $1,536 \cdot 6 = 9,216$ м \\ \hline
		ЛХБ 80 & 5040  & $\frac{44830}{2 \cdot 5040} \approx 5$ & $1,536 \cdot 5 = 7,68$ м \\ \hline
		ЛТБ 80 & 5200  & $\frac{44830}{2 \cdot 5020} \approx 5$ & $1,536 \cdot 5 = 7,68$ м \\ \hline
		ЛБ 80 & 5400  & $\frac{44830}{2 \cdot 5400} \approx 5$ & $1,536 \cdot 5 = 7,68$ м \\ \hline
	\end{tabular}
\end{table}

Сопоставим рассчитанную длину одного ряда светильников $L_\text{р}$ при использовании разных типов люминесцентных ламп с длиной помещения А.

Длина одного ряда светильников ЛОУ 1П при использовании в них ламп типа ЛДЦ 40 больше длины помещения. Следовательно, их применять в осветительной установке в рассматриваемом примере нельзя.

При использовании остальных типов ламп длина одного ряда меньше длины помещения. Целесообразней выбрать вариант с применением ламп ЛД 40, при котором светильники устанавливаются в ряду без разрывов, обеспечивая наиболее равномерное освещение помещения и рабочих мест.

Определяем расчётную освещенность рабочей поверхности Eр в помещении лаборатории при использовании в светильнике ЛОУ 1П двух ламп типа ЛД 40 и относительную погрешность расчёта $\delta_\text{отн}$:

\[
E_\text{р} = 2 \cdot E \cdot \frac{\text{Ф}_\text{л} \cdot n_A}{\text{Ф}_\text{р}} = \frac{2 \cdot 300 \cdot 2500 \cdot 9}{44830} = 301,14 \text{лк}
\]

\[
\delta_\text{отн} = \frac{E_\text{р}-E}{E} \cdot 100 \% = \frac{301,14-300}{300} \cdot 100\%= 0,38 \%
\]

Расчёт считаем законченным, так как отклонение расчётной освещенности $E_\text{р}$ от нормируемой $E$ не превышает допустимые пределы $+20 ÷ −10 \%$.


По результатам расчёта общего искусственного освещения помещения лаборатории построим схему размещения светильников осветительной установки 

\section{Схема размещения}

\tikzset{grid/.style={gray,very thin,opacity=1}}

\begin{figure}
	\centering
	\begin{tikzpicture}[scale=.7]
	\draw[grid] (0,0) grid (21,8);
	\draw[draw=black] (0,0) rectangle ++ (12,8);
	\draw[draw=black] (13,0) rectangle ++ (8,4);
	\end{tikzpicture}
\end{figure}